\section{USAMO 2023}
        \begin{enumerate}
        	\item (\textit{USAMO 2023}) In an acute triangle $ABC$, let $M$ be the midpoint of $\overline{BC}$. Let $P$ be the foot of the perpendicular from $C$ to $AM$. Suppose the circumcircle of triangle $ABP$ intersects line $BC$ at two distinct points $B$ and $Q$. Let $N$ be the midpoint of $\overline{AQ}$. Prove that $NB=NC$.


 



\textbf{Solution 1}

Let $X$ be the foot from $A$ to $\overline{BC}$. By definition, $\angle AXM = \angle MPC = 90^{\circ}$. Thus, $\triangle AXM \sim \triangle CPM$, and $\triangle BMP \sim \triangle AMQ$. 

 From this, we have $\frac{MP}{MX} = \frac{MC}{MA} = \frac{MP}{MQ}$, as $MC=MB$. Thus, $M$ is also the midpoint of $XQ$.

 Now, $NB = NC$ if $N$ lies on the perpendicular bisector of $\overline{BC}$. As $N$ lies on the perpendicular bisector of $\overline{XQ}$, which is also the perpendicular bisector of $\overline{BC}$ (as $M$ is also the midpoint of $XQ$), we are done. ~ Martin2001~ SomebodyST (minor edits)

 



\textbf{Solution 2}

We are going to use barycentric coordinates on $\triangle ABC$. Let $A=(1,0,0)$, $B=(0,1,0)$, $C=(0,0,1)$, and $a=BC$, $b=CA$, $c=AB$. We have $M=\left(0,\frac{1}{2},\frac{1}{2}\right)$ and $P=(x:1:1)$ so $\overrightarrow{CP}=\left(\frac{x}{x+2},\frac{1}{x+2},\frac{1}{x+2}-1\right)$ and $\overrightarrow{AM}=\left(-1,\frac{1}{2},\frac{1}{2}\right)$. Since $\overleftrightarrow{CP}\perp\overleftrightarrow{AM}$, it follows that
 \begin{align*} a^2\left(\frac{1}{2}\cdot\frac{1}{x+2}+\frac{1}{2}\left(\frac{1}{x+2}-1\right)\right)+b^2\left(\frac{1}{2}\cdot\frac{x}{x+2}-\left(\frac{1}{x+2}-1\right)\right)\\ +c^2\left(\frac{1}{2}\cdot\frac{x}{x+2}-\frac{1}{x+2}\right)=0. \end{align*}
Solving this gives
 \[ x=\frac{2b^2-2c^2}{a^2-3b^2-c^2} \]so
 \[ P=\left(\frac{b^2-c^2}{a^2-2b^2-2c^2},\frac{a^2-3b^2-c^2}{2a^2-4b^2-4c^2},\frac{a^2-3b^2-c^2}{2a^2-4b^2-4c^2}\right). \]The equation for $(ABP)$ is
 \[ -a^2yz-b^2zx-c^2xy+ux+vy+wz=0. \]Plugging in $A$ and $B$ gives $u=v=0$. Plugging in $P$ gives
 \begin{align*} -a^2\left(\frac{a^2-3b^2-c^2}{2a^2-4b^2-4c^2}\right)^2-b^2\cdot\frac{a^2-3b^2-c^2}{2a^2-4b^2-4c^2}\cdot\frac{b^2-c^2}{a^2-2b^2-2c^2}\\ -c^2\cdot\frac{b^2-c^2}{a^2-2b^2-2c^2}\cdot\frac{a^2-3b^2-c^2}{2a^2-4b^2-4c^2}+w\cdot\frac{a^2-3b^2-c^2}{2a^2-4b^2-4c^2}=0 \end{align*}
so
 \[ w=\frac{2b^4-2c^4+a^4-3a^2b^2-a^2c^2}{2a^2-4b^2-4c^2}=\frac{a^2}{2}-\frac{b^2}{2}+\frac{c^2}{2}. \]Now let $Q=(0,t,1-t)$ where
 \begin{align*} -a^2t(1-t)+w(1-t)&=0\\ \implies t&=\frac{w}{a^2} \end{align*}
so $Q=\left(0,\frac{w}{a^2},1-\frac{w}{a^2}\right)$. It follows that $N=\left(\frac{1}{2},\frac{w}{2a^2},1-\frac{w}{2a^2}\right)$. It suffices to prove that $\overleftrightarrow{ON}\perp\overleftrightarrow{BC}$. Setting $\overrightarrow{O}=0$, we get $\overrightarrow{N}=\left(\frac{1}{2},\frac{w}{2a^2},1-\frac{w}{2a^2}\right)$. Furthermore we have $\overrightarrow{CB}=(0,1,-1)$ so it suffices to prove that
 \begin{align*} a^2\left(-\frac{w}{2a^2}+\frac{1}{2}-\frac{u}{2a^2}\right)+b^2\left(-\frac{1}{2}\right)+c^2\left(\frac{1}{2}\right)=0\\ \implies w=\frac{a^2}{2}-\frac{b^2}{2}+\frac{c^2}{2} \end{align*}
which is valid. $\square$

 



\textbf{Solution 3}

Let $X$ be the foot from $A$ to $BC.$ Note that $\angle AXC=\angle APC,$ so $(AXPC)$ is cyclic. Furthermore, $AP$ is the radical axis of $(ABP)$ and $(AXPC),$ so since $M$ lies on $AP,$ we must have that $BM\cdot MQ=CM\cdot MX,$ so $M$ is the midpoint of $QX.$ Therefore, $MN\parallel AX\perp BC,$ so we are done. ~krithikrokcs

 


	\item (\textit{USAMO 2023}) Let $\mathbb{R}^{+}$ be the set of positive real numbers. Find all functions $f:\mathbb{R}^{+}\rightarrow\mathbb{R}^{+}$ such that, for all $x, y \in \mathbb{R}^{+}$,

 \[f(xy + f(x)) = xf(y) + 2\]


 



\textbf{Solution 1}

Make the following substitutions to the equation:

 1. $(x, 1) \rightarrow f(x + f(x)) = xf(1) + 2$

 2. $(1, x + f(x)) \rightarrow f(x + f(x) + f(1)) = f(x + f(x)) + 2 = xf(1) + 4$

 3. $(x, 1 + \frac{f(1)}{x}) \rightarrow f(x + f(x) + f(1)) = xf\biggl(1 + \frac{f(1)}{x}\biggr) + 2$

 It then follows from (2) and (3) that $f(1 + \frac{f(1)}{x}) = f(1) + \frac{2}{x}$, so we know that this function is linear for $x > 1$. Substitute $f(x) = ax+b$ and solve for $a$ and $b$ in the functional equation; we find that $f(x) = x + 1 \forall x > 1$.

 Now, we can let $x > 1$ and $y \le 1$. Since $f(x) = x + 1$, $xy + f(x) > x > 1$, so $f(xy + f(x)) = xy + x + 2 = xf(y) + 2$. It becomes clear then that $f(y) = y + 1$ as well, so $f(x) = x + 1$ is the only solution to the functional equation.

 


	\item (\textit{USAMO 2023}) Consider an $n$-by-$n$ board of unit squares for some odd positive integer $n$. We say that a collection $C$ of identical dominoes is a maximal grid-aligned configuration on the board if $C$ consists of $(n^2-1)/2$ dominoes where each domino covers exactly two neighboring squares and the dominoes don't overlap: $C$ then covers all but one square on the board. We are allowed to slide (but not rotate) a domino on the board to cover the uncovered square, resulting in a new maximal grid-aligned configuration with another square uncovered. Let $k(C)$ be the number of distinct maximal grid-aligned configurations obtainable from $C$ by repeatedly sliding dominoes. Find all possible values of $k(C)$ as a function of $n$.


 



\textbf{Solution}

this is solution to jmo3. someone add a solution to amo3

 


 We claim the answer is $(\frac{n+1}{2})^2$. 

 First, consider a checkerboard tiling of the board with 4 colors: R, G, B, Y. Number each column from $1$ to $n$ from left to right and each row from $1$ to $n$ from top to bottom. We color a tile R if its row and column are odd, a tile G is its row is even but its column is odd, a tile B if its row and column is even, and a tile Y if its row is odd but its column is even. 

 Lemma 1: Throughout our moves, the color of the uncolored tile stays an invariant. 

 Consider that a domino can either only change rows or can only change columns. Therefore, sliding a domino into the hole and creating a new one has two possible colors. Of these, note that the new hole will always trivially be two tiles away from the old hole, meaning that the parity of both the row and column number stays the same. Thus, the lemma holds.

 Lemma 2: There are more red tiles than any other color. Because each color is uniquely defined by the parity of a pair of column and row number, it satisfies to show that given an odd integer $n$, there are more odd positive integers less than or equal to $n$ than even ones. Obviously, this is true, and so red will have more tiles than any other color.

 Lemma 3: For any starting configuration $C$ and any blank tile $B$ such that the blank tile's color matches the blank tile's color of $C$, there is no more than one unique configuration $C'$ that can be produced from $C$ using valid moves. 

 We will use proof by contradiction. Assume there exists two different $C'$. We can get from one of these $C'$ to another using moves. However, we have to finish off with the same hole as before. Before the last move, the hole must be two tiles away from the starting hole. However, because the domino we used to move into the blank tile's spot is in the way, that hole must be congruent to the hole produced after the first move. We can induct this logic, and because there is a limited amount of tiles with the same color, eventually we will run out of tiles to apply this to. Therefore, having two distinct $C'$ with the same starting hole $B$ is impossible with some $C$.  

 We will now prove that $(\frac{n+1}{2})^2$ is the answer. There are $\frac{n+1}{2}$ rows and $\frac{n+1}{2}$ columns that are odd, and thus there are $(\frac{n+1}{2})^2$ red tiles. Given lemma 3, this is our upper bound for a maximum. To establish that $(\frac{n+1}{2})^2$ is indeed possible, we construct such a $C$:

 In the first column, leave the first tile up blank. Then, continuously fill in  vertically oriented dominos in that column until it reaches the bottom.In the next $n-1$ columns, place $\frac{n-1}{2}$ vertically oriented dominos in a row starting from the top. At the bottom row, starting with the first unfilled tile on the left, place horizontally aligned dominos in a row until you reach the right.

 Obviously, the top left tile is red. It suffices to show that any red tile may be uncovered. For the first column, one may slide some dominos on the first column until the desired tile is uncovered. For the bottom row, all the first dominos may be slid up, and then the bottom dominos may be slid to the left until the desired red tile is uncovered. Finally, for the rest of the red tiles, the bottom red tile in the same color may be revealed, and then vertically aligned dominos in the same column may be slid down until the desired tile is revealed. Therefore, this configuration may produce $(\frac{n+1}{2})^2$ different configurations with moves. 

 Hence, we have proved that $(\frac{n+1}{2})^2$ is the maximum, and we are done. $\blacksquare{}$

 Do not change the question yourself, SigmaPiE. The question was not maximum, it was all possible values. I changed it back to all possible values. The answer is $\{1,2,...,(\frac{n-1}{2})^2\}\cup\{(\frac{n+1}{2})^2\}$.

 


	\item (\textit{USAMO 2023}) A positive integer $a$ is selected, and some positive integers are written on a board. Alice and Bob play the following game. On Alice's turn, she must replace some integer $n$ on the board with $n+a$, and on Bob's turn he must replace some even integer $n$ on the board with $n/2$. Alice goes first and they alternate turns. If on his turn Bob has no valid moves, the game ends.


 After analyzing the integers on the board, Bob realizes that, regardless of what moves Alice makes, he will be able to force the game to end eventually. Show that, in fact, for this value of $a$ and these integers on the board, the game is guaranteed to end regardless of Alice's or Bob's moves.


 



\textbf{Solution 1}

The contrapositive of the claim is somewhat easier to conceptualize: If it is not guaranteed that the game will end (i.e. the game could potentially last forever), then Bob is not able to force the game to end (i.e. Alice can force it to last forever). So we want to prove that, if the game can potentially last indefinitely, then Alice can force it to last indefinitely.Clearly, if there is $1$ number on the board initially, all moves are forced. This means the claim is true in this specific case, because if the game "potentially" can last forever, this means it must last forever, since the game can only be played in one way. Ergo Alice can "force" this to occur because it is guaranteed to occur. Now we look at all cases where there is more than $1$ number on the board.

 
Case 1: $v_2 (a)=0$

 The game lasts forever here no matter what. This is true because, if the game ends, it means the board was in some position $P$, Alice added $a$ to some number on the board, and all the numbers now on the board are odd. If there are only odd numbers on the board in position $P$, Alice will add $a$ to some number on the board, making it even, meaning this cannot have been the state $P$ of the board. If at least one number n on the board is even, Alice can add $a$ to a number other than $n$, meaning there is still at least 1 even number on the board, meaning this also cannot have been the state $P$ of the board. This covers all possible boards when  $v_2 (a)=0$, so we're done.

 
Case 2: $v_2 (a)=1$

 If there is at least one number n on the board that is even, the game can also last forever. On any move, Alice will add $a$ to this number $n$ if and only if $v_2 (n)=1$. This way, the new number $n'= n+a$ satisfies $v_2 (n') \geq 2$. If Bob does divides $n'$ until $v_2 (n)=1$, Alice will again add $a$ to $n$ resulting in $v_2 (n') \geq 2$. This means that Alice can always keep an even number on the board for Bob to divide no matter how Bob plays.

 If there is no even number on the board, then the game can clearly not last forever. No matter what Alice does, Bob will have no even number to divide after her move.

 
General Case: $v_2 (a)=x$

 In general, it seems to be the case that the game can last indefinitely for some given $a$ if and only if there exists some number $n$ on the board such that $v_2 (n) \geq v_2 (a)$, so we should aim to prove this.

 1. "If"

 Alice can apply a similar strategy to the strategy used in case 2. On any move, Alice will add $a$ to $n$ if and only if $v_2 (n)=v_2 (a)$. If she does this addition, then $v_2 (n') \geq v_2 (a)+1$, keeping an even number on the board. Even if Bob divides $n'$ until $v_2 (n)=v_2 (a)$, Alice will apply the same strategy and keep $v_2 (n') \geq v_2 (a)+1$. Alice's use of this strategy ensures that there always exists some number n on the board such that $v_2 (n) \geq v_2 (a)$, ensuring there always exists an even number n on the board.

 2."Only If"

 If $v_2 (n) < v_2 (a)$ for all n on the board, this means that Alice can never change the value of $v_2 (n)$ for any $n$ on the board. Only Bob can do this, and Bob will subtract $1$ from each $v_2 (n)$ until they are all equal to $0$ (all odd), ending the game.

 We've shown that the game can last indefinitely iff there exists some number $n$ on the board such that $v_2 (n) \geq v_2 (a)$, and have shown that Alice can ensure the game lasts forever in these scenarios using the above strategy. This proves the contrapositive, proving the claim.

 


	\item (\textit{USAMO 2023}) Let $n\geq3$ be an integer. We say that an arrangement of the numbers $1$, $2$, $\dots$, $n^2$ in a $n \times n$ table is row-valid if the numbers in each row can be permuted to form an arithmetic progression, and column-valid if the numbers in each column can be permuted to form an arithmetic progression. For what values of $n$ is it possible to transform any row-valid arrangement into a column-valid arrangement by permuting the numbers in each row?


 



\textbf{Solution 1}

The answer is all $\boxed{\text{prime } n}$.

 \textbf{Proof that primes work} Suppose $n=p$ is prime. Then, let the arithmetic progressions in the $i$th row have least term $a_i$ and common difference $d_i$. For each cell with integer $k$, assign a monomial $x^k$. The sum of the monomials is 
 \[x(1+x+\ldots+x^{n^2-1}) = \sum_{i=1}^n x^{a_i}(1+x^{d_i}+\ldots+x^{(n-1)d_i}),\] where the LHS is obtained by summing over all cells and the RHS is obtained by summing over all rows. Let $S$ be the set of $p$th roots of unity that are not $1$; then, the LHS of the above equivalence vanishes over $S$ and the RHS is 
 \[\sum_{p \mid d_i} x^{a_i}.\] Reducing the exponents (mod $p$) in the above expression  yields 
 \[f(x) := \sum_{p \mid d_i} x^{a_i \pmod{p}} = 0\] when $x \in S$. Note that $\prod_{s \in S} (x-s)=1+x+\ldots+x^{p-1}$ is a factor of $f(x)$, and as $f$ has degree less than $p$, $f$ is either identically 0 or $f(x)=1+x+\ldots+x^{p-1}$. 

 - If $f$ is identically 0, then $p$ never divides $d_i$. Thus, no two elements in each row are congruent $\pmod{p}$, so all residues are represented in each row. Now we can rearrange the grid so that column $i$ consists of all numbers $i \pmod{p}$, which works. 

 - If $f(x)=1+x+\ldots+x^{p-1}$, then $p$ always divides $d_i$. It is clear that each $d_i$ must be $p$, so each row represents a single residue $\pmod{p}$. Thus, we can rearrange the grid so that column $i$ contains all consecutive numbers from $1 + (i-1)p$ to $ip$, which works. 

 
All in all, any prime $n$ satisfies the hypotheses of the problem.

 \textbf{Proof that composites do not work} Let $n=ab$. Look at the term $a^2b+ab$; we claim it cannot be part of a column that has cells forming an arithmetic sequence after any appropriate rearranging. After such a rearrangement, if the column it is in has common difference $d<ab=n$, then $a^2b+ab-d$ must also be in its column, which is impossible. If the column has difference $d > ab = n$, then no element in the next row can be in its column. If the common difference is $d = ab = n$, then $a^2b + ab - 2d = a^2b - ab$ and $a^2b + ab - d = a^2b$, which are both in the row above it, must both be in the same column, which is impossible. Therefore, the grid is not column-valid after any rearrangement, which completes the proof.

 


	\item (\textit{USAMO 2023}) Let ABC be a triangle with incenter $I$ and excenters $I_a$, $I_b$, $I_c$ opposite $A$, $B$, and $C$, respectively. Given an arbitrary point $D$ on the circumcircle of $\triangle ABC$ that does not lie on any of the lines $II_a$, $I_bI_c$, or $BC$, suppose the circumcircles of $\triangle DII_a$ and $\triangle DI_bI_c$ intersect at two distinct points $D$ and $F$. If $E$ is the intersection of lines $DF$ and $BC$, prove that $\angle BAD = \angle EAC$.


 



\textbf{Solution 1}


                    \begin{figure}[H]
                        \centering
                        \includegraphics[width=0.5\textwidth]{Solution/USAMO_I_2023_6_Solution_Figure_1.png}
                    \end{figure}
                    

 Consider points $G,H,J,K,P,$ and $Q$ such that the intersections of the circumcircle of $\triangle{}ABC$ with the circumcircle of $\triangle{}DII_A$ are $D$ and $G$, the intersections of the circumcircle of $\triangle{}ABC$ with the circumcircle of $\triangle{}DI_BI_C$ are $D$ and $H$, the intersections of the circumcircle of $\triangle{}ABC$ with line $\overline{II_A}$ are $A$ and $J$, the intersections of the circumcircle of $\triangle{}ABC$ with line $\overline{I_BI_C}$ are $A$ and $K$, the intersection of lines $\overline{II_A}$ and $\overline{BC}$ is $P$, and the intersection of lines $\overline{I_BI_C}$ and $\overline{BC}$ is $Q$.

 Since $IBI_AC$ is cyclic, the pairwise radical axes of the circumcircles of $\triangle{}DII_A,\triangle{}ABC,$ and $IBI_AC$ concur. The pairwise radical axes of these circles are $\overline{GD},\overline{II_A},$ and $\overline{BC}$, so $G,P,$ and $D$ are collinear. Similarly, since $BCI_BI_C$ is cyclic, the pairwise radical axes of the cirucmcircles of $\triangle{}DI_BI_C,\triangle{}ABC,$ and $BCI_BI_C$ concur. The pairwise radical axes of these circles are $\overline{HD},\overline{I_BI_C},$ and $\overline{BC}$, so $H,Q,$ and $D$ are collinear. This means that $-1=(Q,P;B,C)\stackrel{D}{=}(H,G;B,C)$, so the tangents to the circumcircle of $\triangle{}ABC$ at $G$ and $H$ intersect on $\overline{BC}$. Let this intersection be $X$. Also, let the intersection of the tangents to the circumcircle of $\triangle{}ABC$ at $K$ and $J$ be a point at infinity on $\overline{BC}$ called $Y$ and let the intersection of lines $\overline{KG}$ and $\overline{}HJ$ be $Z$. Then, let the intersection of lines $\overline{GJ}$ and $\overline{HK}$ be $E'$. By Pascal's Theorem on $GGJHHK$ and $GJJHKK$, we get that $X,E',$ and $Z$ are collinear and that $E',Y,$ and $Z$ are collinear, so $E',X,$ and $Y$ are collinear, meaning that $E'$ lies on $\overline{BC}$ since both $X$ and $Y$ lie on $\overline{BC}$.

 
                    \begin{figure}[H]
                        \centering
                        \includegraphics[width=0.5\textwidth]{Solution/USAMO_I_2023_6_Solution_Figure_2.png}
                    \end{figure}
                    

 Consider the transformation which is the composition of an inversion centered at $A$ and a reflection over the angle bisector of $\angle{}CAB$ that sends $B$ to $C$ and $C$ to $B$. We claim that this sends $D$ to $E'$ and $E'$ to $D$. It is sufficient to prove that if the transformation sends $G$ to $G'$, then $AE'JG'$ is cyclic. Notice that $\triangle{}AGB\sim\triangle{}ACG'$ since $\angle{}GAB=\angle{}G'AC$ and $\frac{AG'}{AC}=\frac{\frac{AB\cdot{}AC}{AG}}{AC}=\frac{AB}{AG}$. Therefore, we get that $\angle{}AG'E'=\angle{}ABG=\angle{}AJE'$, so $AE'JG'$ is cyclic, proving the claim. This means that $\angle{}BAE'=\angle{}CAD$.

 
                    \begin{figure}[H]
                        \centering
                        \includegraphics[width=0.5\textwidth]{Solution/USAMO_I_2023_6_Solution_Figure_3.png}
                    \end{figure}
                    

 We claim that $\angle{}I_BE'I_C+\angle{}I_BDI_C=180^\circ$. Construct $D'$ to be the intersection of line $\overline{AE'}$ and the circumcircle of $\triangle{}E'I_BI_C$ and let $B'$ and $C'$ be the intersections of lines $\overline{AC}$ and $\overline{AB}$ with the circumcircle of $\triangle{}BI_BI_C$. Since $B'$ and $C'$ are the reflections of $B$ and $C$ over $\overline{I_BI_C}$, it is sufficient to prove that $A,B',C',D'$ are concyclic. Since $\overline{B'C},\overline{D'E'},$ and $\overline{I_BI_C}$ concur and $D',E',I_B,I_C$ and $I_B,I_C,B',C$ are concyclic, we have that $B',C,D',E'$ are concyclic, so $\angle{}B'D'A=\angle{}ACE'=\angle{}AC'B'$, so $A,B',C',D'$ are concyclic, proving the claim. We can similarly get that $\angle{}IE'I_A=\angle{}IDI_A$.

 
                    \begin{figure}[H]
                        \centering
                        \includegraphics[width=0.5\textwidth]{Solution/USAMO_I_2023_6_Solution_Figure_4.png}
                    \end{figure}
                    

 Let line $\overline{E'J}$ intersect the circumcircle of $\triangle{}DII_A$ at $G$ and $J'$. Notice that $J$ is the midpoint of $\overline{II_A}$ and $\angle{}IE'I_A=\angle{}IDI_A=\angle{}IJ'I_A$, so $IE'I_AJ'$ is a parallelogram with center $J$, so $\frac{EJ}{EJ'}=\frac{1}{2}$. Similarly, we get that if line $\overline{E'K}$ intersects the circumcircle of $\triangle{}DI_BI_C$ at $H$ and $K'$, we have that $\frac{EK}{EK'}=\frac{1}{2}$, so $\overline{KJ}\parallel\overline{K'J'}$, so $\angle{}HGJ'=\angle{}HGJ=\angle{}HKJ=\angle{}HK'J'$, so $G,H,J',K'$ are concyclic. Then, the pairwise radical axes of the circumcircles of $\triangle{}DII_A,\triangle{}DI_BI_C,$ and $GHJ'K'$ are $\overline{DF},\overline{HK'},$ and $\overline{GJ'}$, so $\overline{DF},\overline{HK'},$ and $\overline{GJ'}$ concur, so $\overline{DF},\overline{HK},$ and $\overline{GJ}$ concur, so $E=E'$. We are then done since $\angle{}BAE'=\angle{}CAD$.

 



\textbf{Solution 2}

Set $\triangle ABC$ as the reference triangle, and let $D=(x_0,y_0,z_0)$ with homogenized coordinates. To find the equation of circle $DII_A$, we note that $I=(a:b:c)$ (not homogenized) and $I_A=(-a:b:c)$. Thus, for this circle with equation $a^2yz+b^2xz+c^2xy=(x+y+z)(u_1x+v_1y+w_1z)$, we compute that 
 \[u_1=bc,\] 
 \[v_1=\frac{bc^2x_0}{bz_0-cy_0},\] 
 \[w_1=\frac{b^2cx_0}{cy_0-bz_0}.\]For circle $DI_BI_C$ with equation $a^2yz+b^2xz+c^2xy=(x+y+z)(u_2x+v_2y+w_2z)$,, we find that 
 \[u_2=-bc,\] 
 \[v_2=\frac{bc^2x_0}{bz_0+cy_0},\] 
 \[w_2=\frac{b^2cx_0}{cy_0+bz_0}.\]The equation of the radical axis is $ux+vy+wz=0$ with $u=u_1-u_2$, $v=v_1-v_2$, and $w=w_1-w_2$. We want to consider the intersection of this line with line $\overline{BC}$, so set $x=0$. The equation reduces to $vy+wz=0$. We see that $v=\frac{2bc^3x_0y_0}{b^2z_0^2-c^2y_0^2}$ and $w=\frac{2b^3cx_0z_0}{c^2y_0^2-b^2z_0^2}$, so 
 \[\frac{v}{w}=\frac{b^2z_0}{c^2y_0},\] which is the required condition for $\overline{AD}$ and $\overline{AE}$ to be isogonal.

 


\end{enumerate}